\chapter{Computational Chemistry: Hartree--Fock and DFT}\label{ch:b3:computational-chemistry}

The helium hydride ion, \ce{HeH+}, is thought to be the first molecule that
formed as the young Universe cooled: a helium atom holding a proton. It was
made in a laboratory discharge tube in 1925, but searched for in space for
decades without success, until its rotational line was detected in a
planetary nebula in 2019. By then its bond length, its vibrational
frequencies and the wavelengths of its lines had long been known --- from
calculation. A computer program that solves, approximately, the
\omterm{thm:b3:quantum-model-systems:schrodinger}{Schrödinger equation} of the electrons can predict the shape, the energy
and the spectrum of a molecule nobody has put in a bottle. This chapter
explains how such programs work, from the separation of nuclei and electrons
to Hartree--Fock theory and \omterm{def:b3:computational-chemistry:dft}{density-functional theory}, and what their numbers
can and cannot be trusted for.

\begin{recall}
\Cref{ch:b3:quantum-model-systems} gave the Hamiltonian, the \omterm{thm:b3:quantum-model-systems:schrodinger}{Schrödinger
equation} and the hydrogen atom; \cref{ch:b3:many-electron-atoms} the \omterm{def:b3:many-electron-atoms:slater}{Slater
determinant} and the \omterm{def:b3:many-electron-atoms:exchange}{exchange integral}. The Year 2 volume built molecular
orbitals as linear combinations of atomic orbitals (LCAO), with overlap,
Coulomb and resonance integrals and a secular determinant, and solved the
Hückel method; it also used photoelectron spectra to see orbital energies.
\end{recall}

\begin{omfigure}
\includegraphics[width=0.6\linewidth]{images/book4/ai/b3-03-cluster.jpg}

{\small A computing cluster: rows of servers on which most quantum-chemical
calculations now run, often for hours or days per molecule.}
\end{omfigure}

\section{Born--Oppenheimer and the potential energy surface}

A molecule's Hamiltonian contains the kinetic energies of its nuclei and
electrons and all their Coulomb interactions. Nuclei are at least 1836
times heavier than electrons: for the same forces they move far more
slowly, and the electrons adjust almost instantly to each position of the
nuclei.

\begin{definition}[Born--Oppenheimer approximation]\label{def:b3:computational-chemistry:born-oppenheimer}
In the \emph{Born--Oppenheimer approximation}\index{Born--Oppenheimer
approximation}, the electronic \omterm{thm:b3:quantum-model-systems:schrodinger}{Schrödinger equation} is solved for nuclei
held fixed at positions $\mathbf R$,
\[
\hat H_{\mathrm{el}}\,\psi_{\mathrm{el}}(\mathbf r;\mathbf R) =
E_{\mathrm{el}}(\mathbf R)\,\psi_{\mathrm{el}}(\mathbf r;\mathbf R),
\]
and the nuclei then move in the potential $E_{\mathrm{el}}(\mathbf R) +
V_{\mathrm{nn}}(\mathbf R)$, the electronic energy plus the repulsion of the
nuclei.
\end{definition}

\begin{definition}[Potential energy surface, stationary point, geometry optimisation]\label{def:b3:computational-chemistry:pes}
The \emph{potential energy surface}\index{potential energy surface} (PES) of
a molecule is the function $U(\mathbf R) = E_{\mathrm{el}}(\mathbf R) +
V_{\mathrm{nn}}(\mathbf R)$ of its nuclear coordinates. A \emph{stationary
point}\index{stationary point} is a point where all first derivatives of $U$
vanish; a minimum is an \emph{equilibrium geometry}\index{equilibrium
geometry}. Searching for one from a starting structure, by following the
forces $-\nabla U$, is a \emph{geometry optimisation}\index{geometry
optimisation}.
\end{definition}

For a diatomic molecule the PES is a curve $U(R)$; for a triatomic a
function of three coordinates; for $N$ atoms of $3N - 6$. The mass of the
nuclei does not enter $\hat H_{\mathrm{el}}$: \ce{H2}, \ce{HD} and \ce{D2} share the
same PES, the same bond length and \omterm{def:b3:quantum-model-systems:oscillator}{force constant}, and differ only in how
the nuclei move on it (\omterm{def:b3:quantum-model-systems:oscillator}{zero-point energy}, \cref{ch:b3:quantum-model-systems}).

\begin{proposition}[Frequencies from the Hessian]\label{prop:b3:computational-chemistry:hessian}
Near a minimum, $U \approx U_0 + \frac12\sum_{ij}H_{ij}\,\delta q_i\,\delta q_j$, with
the Hessian $H_{ij} = \partial^2U/\partial q_i\partial q_j$. In mass-weighted
coordinates $\delta q_i\sqrt{m_i}$, the \omterm{def:b3:quantum-model-systems:operator}{eigenvalues} $\lambda_k$ of the
mass-weighted Hessian give the harmonic angular frequencies $\omega_k =
\sqrt{\lambda_k}$ of the normal modes. At a minimum all $\lambda_k$ are
positive (six, or five for a linear molecule, are zero: translations and
rotations); at a first-order saddle point exactly one is negative.
\end{proposition}

\begin{proof}[Partial proof]
For a diatomic, $U \approx U_0 + \frac12k(R - R_e)^2$ with $k = U''(R_e)$, and
the oscillator of \cref{ch:b3:quantum-model-systems} gives $\omega =
\sqrt{k/\mu}$. In general, the classical equations $m_i\ddot q_i =
-\sum_jH_{ij}q_j$ become, in mass-weighted coordinates $x_i = \sqrt{m_i}q_i$,
$\ddot{\mathbf x} = -\tilde H\mathbf x$ with the symmetric matrix $\tilde H_{ij} =
H_{ij}/\sqrt{m_im_j}$; its eigenvectors oscillate independently at
$\sqrt{\lambda_k}$. A negative $\lambda_k$ gives an imaginary frequency: along
that direction the energy goes down, a maximum. The normal modes of
\cref{ch:b3:group-theory-applied} are these eigenvectors.
\end{proof}

\begin{omfigure}
\begin{tikzpicture}[font=\small, scale=0.95]
  \fill[omDef!8] (0,0) ellipse (1.9 and 1.15);
  \foreach \s in {1.0,0.75,0.5,0.25}{\draw[omDef!70] (0,0) ellipse ({1.9*\s} and {1.15*\s});}
  \fill[omThm!8] (6,0.4) ellipse (1.6 and 1.0);
  \foreach \s in {1.0,0.75,0.5,0.25}{\draw[omThm!70] (6,0.4) ellipse ({1.6*\s} and {1.0*\s});}
  \draw[gray] (2.1,-0.15) .. controls (2.9,0.55) and (3.4,0.55) .. (4.3,0.35);
  \draw[gray] (2.2,0.9) .. controls (3.0,0.35) and (3.5,0.35) .. (4.1,1.2);
  \node[omsaddle] at (3.15,0.42) {};
  \draw[ompath] (0,0) .. controls (1.8,0.6) and (2.4,0.5) .. (3.15,0.42) .. controls (4.0,0.35) and (4.8,0.4) .. (6,0.4);
  \node[anchor=north] at (0,-1.25) {minimum A (reactants)};
  \node[anchor=north] at (6,-0.7) {minimum B (products)};
  \node[anchor=south] at (3.15,1.25) {saddle point};
  \draw[->, gray] (3.15,1.25) -- (3.15,0.6);
  \fill (0,0) circle (1.2pt) (6,0.4) circle (1.2pt);
  \draw[->] (-2.3,-2.2) -- (8.2,-2.2) node[right] {$q_1$};
  \draw[->] (-2.3,-2.2) -- (-2.3,1.8) node[above] {$q_2$};
\end{tikzpicture}

{\small A \omterm{def:b3:computational-chemistry:pes}{potential energy surface} seen from above, as a contour map over two
nuclear coordinates (schematic). Two minima are joined, through a saddle
point (cross), by the path of lowest energy (dashed): the minimum energy path
studied in \cref{ch:b3:rate-theories}.}
\end{omfigure}

\section{The variational principle}

The electronic \omterm{thm:b3:quantum-model-systems:schrodinger}{Schrödinger equation} cannot be solved exactly for any
molecule with more than one electron. Approximations are judged by one
theorem.

\begin{theorem}[The variational principle]\label{thm:b3:computational-chemistry:variational}
For any normalisable trial function $\phi$ satisfying the boundary
conditions of the problem,
\[
E[\phi] = \frac{\langle\phi|\hat H|\phi\rangle}{\langle\phi|\phi\rangle} \ge E_0,
\]
where $E_0$ is the lowest \omterm{def:b3:quantum-model-systems:operator}{eigenvalue} of $\hat H$; equality holds only if $\phi$
is a ground-state \omterm{def:b3:quantum-model-systems:operator}{eigenfunction}. This is the \emph{variational
principle}\index{variational principle}.
\end{theorem}

\begin{proof}
Expand $\phi = \sum_nc_n\psi_n$ on the orthonormal \omterm{def:b3:quantum-model-systems:operator}{eigenfunctions} of $\hat H$
(\cref{thm:b3:quantum-model-systems:hermitian-real}). Then $\langle\phi|\hat
H|\phi\rangle = \sum_n|c_n|^2E_n \ge E_0\sum_n|c_n|^2 = E_0\langle\phi|\phi\rangle$,
since every $E_n \ge E_0$; equality requires $c_n = 0$ whenever $E_n > E_0$.
\end{proof}

The lower the energy of a trial function, the better it is; any adjustable
parameter is fixed by minimising $E$.

\begin{proposition}[Helium with a screened charge]\label{prop:b3:computational-chemistry:helium-zeta}
For helium, the trial function $\phi = \eu^{-\zeta(r_1 + r_2)}$ (in atomic units:
lengths in $a_0$, energies in $E_h$) gives $E(\zeta) = \zeta^2 - \frac{27}{8}\zeta$,
minimal for $\zeta = \frac{27}{16}$, with $E = -(\frac{27}{16})^2 E_h =
\qty{-2.8477}{\hartree}$.
\end{proposition}

\begin{proof}
For one electron in $\eu^{-\zeta r}$, $\langle T\rangle = \zeta^2/2$ and
$\langle 1/r\rangle = \zeta$; the repulsion of two electrons in the same $1s$
function is $\frac58\zeta$ (a standard integral, admitted). So $E = 2\cdot\frac12
\zeta^2 - 2\cdot2\zeta + \frac58\zeta = \zeta^2 - \frac{27}{8}\zeta$, and $\dd E/\dd
\zeta = 0$ gives $\zeta = 27/16 = 1.6875$.
\end{proof}

\begin{example}[How good is it?]\label{ex:b3:computational-chemistry:helium}
% ledger: ie:He, ie4:He2, const:Eh, const:eV
The total energy of helium is minus the sum of its two ionisation energies,
$-(24.5874 + 54.4178)\,\unit{eV} = \qty{-2.9034}{\hartree}$. The one-parameter
function is \qty{0.056}{\hartree} (\qty{1.5}{eV}) too high: each electron sees a
nucleus screened to $Z_{\mathrm{eff}} = 1.69$ by the other, as Slater's rules
suggest, but the electrons also dodge each other in a way no product of two
functions can describe.
\end{example}

\begin{proposition}[Linear variation]\label{prop:b3:computational-chemistry:linear-variation}
For a trial function $\phi = \sum_{i=1}^n c_i\chi_i$ built on $n$ fixed functions,
with $H_{ij} = \langle\chi_i|\hat H|\chi_j\rangle$ and $S_{ij} = \langle\chi_i|\chi_j
\rangle$, the stationary values of $E$ are the roots of $\det(H - ES) = 0$, and
the lowest root is an upper bound to $E_0$.
\end{proposition}

\begin{proof}
$E\sum_{ij}c_ic_jS_{ij} = \sum_{ij}c_ic_jH_{ij}$ (real coefficients). Differentiate
with respect to $c_k$ and set $\partial E/\partial c_k = 0$: $\sum_j(H_{kj} -
ES_{kj})c_j = 0$ for every $k$. A non-zero solution needs a vanishing
determinant. The lowest root is $E[\phi]$ for its eigenvector, hence $\ge E_0$.
\end{proof}

The Hückel method of the Year 2 volume is this proposition with $p$
orbitals and empirical $H_{ij}$. Every method below is the same idea with
better functions and better matrix elements.

\section{Basis sets}

\begin{definition}[Basis set, minimal basis, Slater- and Gaussian-type orbitals]\label{def:b3:computational-chemistry:basis}
A \emph{basis set}\index{basis set} is the set of fixed functions $\chi_i$,
centred on the atoms, from which molecular orbitals are built. A
\emph{minimal basis set}\index{minimal basis set} has one function per
occupied atomic orbital (one $1s$ for H, five for C). A \emph{Slater-type
orbital}\index{Slater-type orbital} (STO) has the radial form
$r^{n-1}\eu^{-\zeta r}$; a \emph{Gaussian-type orbital}\index{Gaussian-type
orbital} (GTO) has the form $\eu^{-\alpha r^2}$ times a polynomial in $x$, $y$,
$z$. A \emph{contracted Gaussian function}\index{contracted Gaussian function}
is a fixed combination $\sum_kd_k\,\eu^{-\alpha_kr^2}$ of primitive Gaussians.
\end{definition}

Slater functions have the right shape (a cusp at the nucleus, an
exponential tail) but their two-electron integrals over several centres are
very costly; Gaussian functions have the wrong shape, but the product of two
Gaussians on different centres is a Gaussian on a point between them, so
every integral reduces to a closed formula. Contractions get the best of
both: in the STO-3G basis, each Slater function is replaced by a fixed sum of
three Gaussians fitted to it.

\begin{example}[The STO-3G function of hydrogen]\label{ex:b3:computational-chemistry:sto3g}
% ledger: bse:H.sto3g.a1, bse:H.sto3g.a2, bse:H.sto3g.a3, bse:sto3g.c1, bse:sto3g.c2, bse:sto3g.c3
The hydrogen $1s$ function of STO-3G is $0.15433\,g(3.42525) + 0.53533\,
g(0.62391) + 0.44463\,g(0.16886)$, where $g(\alpha)$ is the normalised Gaussian
$(2\alpha/\pi)^{3/4}\eu^{-\alpha r^2}$ (exponents in $a_0^{-2}$). It imitates a
Slater function of exponent $\zeta = 1.24$ --- a hydrogen atom slightly
compressed, as it is in molecules. Its overlap with the exact $1s$ function
of the same exponent is $0.9998$; only the cusp at the nucleus and the far
tail are missed (figure below).
\end{example}

\begin{definition}[Split-valence, polarisation and diffuse functions]\label{def:b3:computational-chemistry:split-valence}
A \emph{split-valence basis set}\index{split-valence basis set} describes each
valence orbital by two (or more) functions of different sizes, so that the
orbital can grow or shrink in a molecule. A \emph{polarisation
function}\index{polarisation function} has a higher $l$ than the occupied
orbitals ($p$ on H, $d$ on C) and lets the density shift off the atom; a
\emph{diffuse function}\index{diffuse function} has a small exponent and
describes anions and weak interactions far from the nuclei.
\end{definition}

A name such as 6-31G(d) says: six primitives for each core function, the
valence split into a contraction of three and a single primitive, and $d$
\omterm{def:b3:computational-chemistry:split-valence}{polarisation functions} on heavy atoms. Results converge towards the
\emph{basis-set limit} as the basis grows; a calculation is only as good as
its basis.

\begin{omfigure}
\begin{tikzpicture}
\begin{axis}[width=0.46\linewidth, height=5.4cm, xmin=0, xmax=5, ymin=0,
  ymax=0.6, axis lines=left, xlabel={$r$ / $a_0$}, ylabel={$\phi(r)$},
  legend cell align=left, legend style={font=\small, draw=none},
  tick label style={font=\small}, label style={font=\small},
  title={\small STO-3G against $\eu^{-r}$ ($\zeta = 1$)}]
\addplot[black, thick] table[x=r, y=slater] {figdata/out/bachelor-3-minimal-basis-hf-sto.dat};
\addlegendentry{Slater}
\addplot[omDef, thick, dashed] table[x=r, y=g3] {figdata/out/bachelor-3-minimal-basis-hf-sto.dat};
\addlegendentry{STO-3G}
\addplot[omThm, thick, dotted] table[x=r, y=g1] {figdata/out/bachelor-3-minimal-basis-hf-sto.dat};
\addlegendentry{one Gaussian}
\end{axis}
\end{tikzpicture}\hfill
\begin{tikzpicture}
\begin{axis}[width=0.46\linewidth, height=5.4cm, xmin=0.5, xmax=8.5,
  ymin=-2.86, ymax=-2.78, axis lines=left, xlabel={SCF cycle},
  ylabel={$E$ / $E_h$}, ytick={-2.85,-2.83,-2.81,-2.79},
  yticklabel style={/pgf/number format/fixed, /pgf/number format/precision=2},
  tick label style={font=\small}, label style={font=\small},
  title={\small SCF of \ce{HeH+}}]
\addplot[omDef, thick, mark=*] table[x=it, y=E] {figdata/out/bachelor-3-minimal-basis-hf-heh-scf.dat};
\addplot[gray, dashed, domain=0.5:8.5] {-2.84184};
\end{axis}
\end{tikzpicture}

{\small Left: the STO-3G contraction (dashed) reproduces the $1s$ Slater
function closely except at the nucleus; a single Gaussian (dotted) does not.
Right: the energy of \ce{HeH+} at \qty{1.4632}{\bohr} during the
self-consistent-field iterations (STO-3G); it converges from above to
\qty{-2.8418}{\hartree} within a few cycles, each iterate obeying the \omterm{thm:b3:computational-chemistry:variational}{variational
principle}.}
\end{omfigure}

\section{Hartree--Fock theory}

The simplest \omterm{def:b3:many-electron-atoms:indistinguishable}{antisymmetric wavefunction} of $N$ electrons is one \omterm{def:b3:many-electron-atoms:slater}{Slater
determinant}. Hartree--Fock theory finds the best one.

\begin{definition}[Hartree--Fock method, Fock operator, self-consistent field]\label{def:b3:computational-chemistry:hartree-fock}
The \emph{Hartree--Fock method}\index{Hartree--Fock method} approximates the
ground state by the single \omterm{def:b3:many-electron-atoms:slater}{Slater determinant} of lowest energy. Its orbitals
are \omterm{def:b3:quantum-model-systems:operator}{eigenfunctions} of the \emph{Fock operator}\index{Fock operator}
\[
\hat f(1) = \hat h(1) + \sum_b\big[2\hat J_b(1) - \hat K_b(1)\big],
\]
in the closed-shell case: $\hat h$ is the kinetic energy and attraction to the
nuclei of one electron, $\hat J_b$ the repulsion by the charge cloud of
orbital $b$, and $\hat K_b$ the exchange \omterm{def:b3:quantum-model-systems:operator}{operator}, whose \omterm{def:b3:quantum-model-systems:hermitian}{expectation values} are
the \omterm{def:b3:many-electron-atoms:exchange}{exchange integrals} of \cref{ch:b3:many-electron-atoms}. Since $\hat f$
depends on the orbitals it determines, the equations are solved by
iteration until the orbitals no longer change: a \emph{self-consistent
field}\index{self-consistent field} (SCF).
\end{definition}

\begin{theorem}[Roothaan--Hall equations]\label{thm:b3:computational-chemistry:roothaan}
Writing each orbital in a basis, $\phi_a = \sum_\mu C_{\mu a}\chi_\mu$, the
closed-shell Hartree--Fock equations become the matrix eigenproblem
\[
\mathbf F\mathbf C = \mathbf S\mathbf C\boldsymbol\varepsilon, \qquad F_{\mu\nu} =
h_{\mu\nu} + \sum_{\lambda\sigma}P_{\lambda\sigma}\big[(\mu\nu|\sigma\lambda) -
\tfrac12(\mu\lambda|\sigma\nu)\big],
\]
with the density matrix $P_{\lambda\sigma} = 2\sum_a^{\mathrm{occ}}C_{\lambda a}C_{\sigma
a}$ and the two-electron integrals $(\mu\nu|\lambda\sigma) = \iint\chi_\mu(1)\chi_\nu(1)
r_{12}^{-1}\chi_\lambda(2)\chi_\sigma(2)$ (atomic units).
\end{theorem}

\begin{proof}[Partial proof]
The energy of the determinant is $E = \sum_{\mu\nu}P_{\mu\nu}h_{\mu\nu} + \frac12
\sum P_{\mu\nu}P_{\lambda\sigma}[(\mu\nu|\sigma\lambda) - \frac12(\mu\lambda|\sigma\nu)]$, a
quadratic function of the coefficients. Minimise it under the constraints
that the orbitals stay orthonormal, $\mathbf C^{\mathsf T}\mathbf S\mathbf C = \mathbf 1$,
with a Lagrange multiplier for each constraint: the stationarity
conditions are $\mathbf F\mathbf C = \mathbf S\mathbf C\boldsymbol\varepsilon$, the
multipliers forming a matrix that can be made diagonal by a rotation of the
occupied orbitals among themselves. The algebra is admitted.
\end{proof}

\begin{method}[The SCF procedure]\label{met:b3:computational-chemistry:scf}
\begin{enumerate}
  \item Compute the integrals $S_{\mu\nu}$, $h_{\mu\nu}$ and $(\mu\nu|\lambda\sigma)$ once.
  \item Guess the density matrix (for instance from the eigenvectors of
        $\mathbf h$ alone, the ``core'' guess).
  \item Build $\mathbf F$ from $\mathbf P$; solve $\mathbf F\mathbf C = \mathbf S\mathbf C
        \boldsymbol\varepsilon$ (orthogonalise with $\mathbf S^{-1/2}$, then diagonalise).
  \item Fill the lowest orbitals, two electrons each; form the new $\mathbf P$
        and the energy.
  \item Repeat from step 3 until the energy and $\mathbf P$ change by less than a
        threshold.
\end{enumerate}
\end{method}

\begin{omfigure}
\begin{tikzpicture}[font=\small, node distance=4mm,
  box/.style={draw=black!60, rounded corners=2pt, fill=omDef!6, align=center,
  text width=4.4cm, inner sep=4pt}]
  \node[box] (a) {integrals $S$, $h$, $(\mu\nu|\lambda\sigma)$};
  \node[box, below=of a] (b) {guess $\mathbf P$};
  \node[box, below=of b] (c) {build $\mathbf F(\mathbf P)$};
  \node[box, below=of c] (d) {solve $\mathbf F\mathbf C = \mathbf S\mathbf C\boldsymbol\varepsilon$};
  \node[box, below=of d] (e) {new $\mathbf P$, energy $E$};
  \node[diamond, aspect=2.4, draw=black!60, fill=omThm!6, below=of e, inner sep=1pt] (f) {converged?};
  \node[box, below=of f, fill=omProp!8] (g) {orbitals, $\varepsilon$, $E$};
  \draw[->] (a) -- (b); \draw[->] (b) -- (c); \draw[->] (c) -- (d);
  \draw[->] (d) -- (e); \draw[->] (e) -- (f);
  \draw[->] (f) -- node[left] {yes} (g);
  \draw[->] (f.east) -- ++(1.6,0) |- node[pos=0.25, right] {no} (c.east);
\end{tikzpicture}

{\small The self-consistent-field loop of a Hartree--Fock calculation.}
\end{omfigure}

\begin{example}[\ce{H2} in the minimal basis]\label{ex:b3:computational-chemistry:h2}
% ledger: re:H2, janaf:H.dfH0, diat:H2.we, diat:H2.wexe
With one STO-3G function per atom, the bonding orbital is fixed by symmetry,
$\sigma_g \propto \chi_A + \chi_B$, and the SCF needs a single step. At $R =
\qty{1.4}{\bohr}$ the overlap is $S_{AB} = 0.6593$ and the energy
\qty{-1.1167}{\hartree}; the minimum lies at \qty{1.346}{\bohr} (\qty{71.2}{pm},
against the measured \qty{74.1}{pm}) and \qty{-1.1175}{\hartree}. The figure below
compares the curve with the true one, built from measured spectroscopic
data and the dissociation energy (a Morse curve, \cref{ch:b3:rovibrational-spectroscopy}).
\end{example}

\begin{omfigure}
\begin{tikzpicture}
\begin{axis}[width=0.7\linewidth, height=6cm, xmin=0.8, xmax=10, ymin=-1.2,
  ymax=-0.45, axis lines=left, xlabel={$R$ / $a_0$}, ylabel={$E$ / $E_h$},
  legend cell align=left, legend style={font=\small, at={(0.04,0.98)}, anchor=north west, draw=none},
  tick label style={font=\small}, label style={font=\small}]
\addplot[omDef, thick] table[x=R, y=rhf] {figdata/out/bachelor-3-minimal-basis-hf-h2.dat};
\addlegendentry{RHF/STO-3G}
\addplot[black, thick, dashed] table[x=R, y=morse] {figdata/out/bachelor-3-minimal-basis-hf-h2.dat};
\addlegendentry{measured (Morse form)}
\addplot[gray, dotted, domain=0.8:10] {-1.0};
\addplot[omThm, dotted, domain=0.8:10] {-0.93316};
\node[font=\scriptsize, anchor=north] at (axis cs:8,-1.005) {two exact H atoms};
\node[font=\scriptsize, anchor=south, omThm] at (axis cs:8,-0.930) {two STO-3G H atoms};
\end{axis}
\end{tikzpicture}

{\small The energy of \ce{H2} against the internuclear distance: restricted
Hartree--Fock in the STO-3G basis (solid) and the true curve (dashed, Morse
form built from the measured $D_0$, $\omega_e$, $\omega_ex_e$, $r_e$). Near the
minimum RHF is good; at large $R$ it rises far above two hydrogen atoms.}
\end{omfigure}

\begin{proposition}[Why restricted Hartree--Fock fails at dissociation]\label{prop:b3:computational-chemistry:rhf-dissociation}
In the minimal basis, the RHF wavefunction of \ce{H2} is, apart from
normalisation, $\sigma_g(1)\sigma_g(2) \propto \chi_A(1)\chi_B(2) + \chi_B(1)\chi_A(2)
+ \chi_A(1)\chi_A(2) + \chi_B(1)\chi_B(2)$: at any distance, half covalent and half
ionic (\ce{H+ H-}). At large $R$ its energy tends to the mean of two
neutral atoms and of an ion pair, not to that of two atoms.
\end{proposition}

\begin{proof}
Expand $(\chi_A + \chi_B)(1)(\chi_A + \chi_B)(2)$: the two cross terms put one
electron on each atom, the two square terms both on the same atom, with
equal weights whatever $R$. The determinant cannot change these weights.
\end{proof}

\begin{definition}[Electron correlation, correlation energy]\label{def:b3:computational-chemistry:correlation}
\emph{Electron correlation}\index{electron correlation} is the part of the
electrons' mutual avoidance not described by a single determinant (whose
electrons of opposite spins move independently in each other's average
field). The \emph{correlation energy}\index{correlation energy} is the exact
energy minus the Hartree--Fock energy in a complete basis; it is always
negative.
\end{definition}

The \omterm{def:b3:computational-chemistry:correlation}{correlation energy} is about \qty{-0.04}{\hartree} per electron pair, one
per cent of a total energy but comparable to the energy of a reaction:
methods that add it (configuration interaction, perturbation theory,
coupled clusters) cost far more than Hartree--Fock, whose cost already grows
roughly as the fourth power of the number of basis functions.

\begin{theorem}[Koopmans' theorem]\label{thm:b3:computational-chemistry:koopmans}
In Hartree--Fock theory, the energy needed to remove an electron from
occupied orbital $a$, keeping all other orbitals frozen, is $-\varepsilon_a$. The
measured ionisation energies are therefore approximately the negatives of
the orbital energies: this is \emph{Koopmans' theorem}\index{Koopmans'
theorem}.
\end{theorem}

\begin{proof}
The energy of a closed-shell determinant is $\sum_a2h_{aa} + \sum_{ab}(2J_{ab} -
K_{ab})$, and $\varepsilon_a = h_{aa} + \sum_b(2J_{ab} - K_{ab})$. Remove one electron
from $a$ without changing the orbitals: the lost terms are $h_{aa}$, the
interactions of that electron with all others, $\sum_b(2J_{ab} - K_{ab}) -
J_{aa}$, plus $J_{aa}$ with its former partner --- in total exactly
$\varepsilon_a$. So $E^+ - E = -\varepsilon_a$.
\end{proof}

\begin{example}[Ionising \ce{H2}]\label{ex:b3:computational-chemistry:koopmans}
% ledger: ie4:H2
At \qty{1.4}{\bohr}, the STO-3G orbital energy of $\sigma_g$ is
\qty{-0.5782}{\hartree}: Koopmans predicts \qty{15.73}{eV}, against the measured
\qty{15.43}{eV}. The frozen orbitals overestimate the energy of the ion (it
would relax), the missing correlation acts the other way; for valence
ionisations the two errors largely cancel, which is why the photoelectron
spectra of the Year 2 volume can be read with orbital diagrams.
\end{example}

\section{Density-functional theory, and what a calculation can tell}

The wavefunction of $N$ electrons depends on $3N$ coordinates; the
\omterm{def:b3:computational-chemistry:dft}{electron density} depends on three.

\begin{definition}[Electron density, density-functional theory]\label{def:b3:computational-chemistry:dft}
The \emph{electron density}\index{electron density} $\rho(\mathbf r)$ is the
number of electrons per unit volume at $\mathbf r$, summed over all electrons:
$\int\rho\,\dd\tau = N$. \emph{Density-functional theory}\index{density-functional
theory} (DFT) computes the ground-state energy as a functional of $\rho$. In
practice $\rho = \sum_a|\phi_a|^2$ is built from \emph{Kohn--Sham
orbitals}\index{Kohn--Sham orbitals}, the orbitals of fictitious independent
electrons that have the same density as the real ones, and all the effects
of exchange and correlation are gathered in an \emph{exchange--correlation
functional}\index{exchange--correlation functional} $E_{\mathrm{xc}}[\rho]$.
\end{definition}

\begin{theorem}[Hohenberg--Kohn]\label{thm:b3:computational-chemistry:hohenberg-kohn}
The ground-state density of a system of electrons determines the external
potential (the nuclei) up to a constant, hence the Hamiltonian and every
ground-state property; and the energy functional of the density is
minimal at the true density.
\end{theorem}
\admitted

The proof is treated in more advanced courses. It guarantees that an exact
functional exists, not what it is: every practical functional is an
approximation, chosen from a hierarchy --- local density, gradient
corrections, hybrids mixing in some Hartree--Fock exchange --- and
calibrated on reference data. The Kohn--Sham equations are solved by the
same SCF loop as Hartree--Fock, at a similar cost, but include correlation;
this is why DFT became the workhorse of chemistry.

\begin{method}[Reading a calculation]\label{met:b3:computational-chemistry:read-calculation}
\begin{enumerate}
  \item Check that the SCF converged and that the \omterm{def:b3:computational-chemistry:pes}{geometry optimisation} ended
        with forces below threshold.
  \item Run a frequency calculation at the same level: all real frequencies
        mean a minimum; one imaginary frequency, a transition structure, whose
        normal mode is the motion across the barrier.
  \item Compare like with like: energy differences between structures
        computed at one level, never absolute energies from different levels.
  \item Compare with experiment knowing the typical errors of the level:
        harmonic frequencies are usually scaled down by 0.9 to 0.97;
        bond lengths are good to about 1 pm with a hybrid functional and a
        polarised basis; reaction energies to a few \unit{kJ/mol} at best.
\end{enumerate}
\end{method}

\begin{method}[Choosing a level of theory]\label{met:b3:computational-chemistry:choose-level}
Geometries and frequencies of organic molecules: a hybrid functional with a
polarised split-valence basis. Weak interactions: add a dispersion
correction and \omterm{def:b3:computational-chemistry:split-valence}{diffuse functions}. Reaction barriers and accurate
energetics: correlated wavefunction methods on DFT geometries, as far as the
budget allows. Large systems (proteins, surfaces): a small region treated
quantum-mechanically inside a classical force field. Always test the chosen
level on a related molecule whose answer is known.
\end{method}

\begin{inthelab}[A computational experiment]
A calculation is run like an experiment and recorded in the notebook: the
program and its version, the method, the basis, the convergence thresholds,
the starting structure, and the output files kept as raw data. A typical
sequence is: build the molecule, optimise its geometry at a modest level,
compute frequencies to confirm a minimum, then compute the energy at a
higher level on that geometry.
\end{inthelab}

\begin{history}[A molecule known first by calculation]
% ledger: hist4:HeHp.2019
\ce{HeH+} was produced in the laboratory in 1925. Its properties were computed
in detail from the 1930s onwards, and it was predicted to be abundant in
some astrophysical gases; but its first rotational line lies in the far
infrared, absorbed by the atmosphere. In 2019 an airborne observatory
detected it in the planetary nebula NGC 7027, at the frequency calculations
and laboratory spectra had predicted.
\end{history}

\section{Exercises}

\begin{exercise}[$\star$]\label{exo:b3:computational-chemistry:1}
For the hydrogen atom, take the trial function $\eu^{-\alpha r^2}$. In atomic
units $E(\alpha) = \frac32\alpha - 2\sqrt{2\alpha/\pi}$. Find the best $\alpha$ and
the corresponding energy; compare with $-\frac12E_h$.
\end{exercise}

\begin{exercise}[$\star$]\label{exo:b3:computational-chemistry:2}
How many basis functions does a calculation on benzene use in the STO-3G
basis? In 6-31G(d), with six Cartesian $d$ functions per carbon and two $s$
functions per hydrogen?
\end{exercise}

\begin{exercise}[$\star$]\label{exo:b3:computational-chemistry:3}
Explain why \ce{H2} and \ce{D2} have the same bond length and the same \omterm{def:b3:quantum-model-systems:oscillator}{force
constant} but different vibrational wavenumbers. Which ratio do you expect?
\end{exercise}

\begin{exercise}[$\star$]\label{exo:b3:computational-chemistry:4}
A frequency calculation on a structure of formula \ce{C2H5F} gives 17 real
frequencies and one imaginary frequency, $487\iu$~\unit{cm^{-1}}. What kind of
\omterm{def:b3:computational-chemistry:pes}{stationary point} is it? How many frequencies did you expect in total?
\end{exercise}

\begin{exercise}[$\star\star$]\label{exo:b3:computational-chemistry:5}
% ledger: ie:He, ie4:He2
Using $E(\zeta) = \zeta^2 - \frac{27}{8}\zeta$, compute the best variational
energy of helium in \unit{\hartree} and eV, and its error against the exact
$\qty{-2.9034}{\hartree}$. What fraction of the total energy is the error?
\end{exercise}

\begin{exercise}[$\star\star$]\label{exo:b3:computational-chemistry:6}
Two basis functions give $H_{11} = \qty{-1.0}{\hartree}$, $H_{22} =
\qty{-0.5}{\hartree}$, $H_{12} = \qty{-0.2}{\hartree}$, $S_{12} = 0.3$. Solve the
secular equation. Is the lower root below $H_{11}$?
\end{exercise}

\begin{exercise}[$\star\star$]\label{exo:b3:computational-chemistry:7}
A Hartree--Fock calculation takes 10 minutes for benzene in STO-3G. Assuming
a cost proportional to the fourth power of the number of basis functions,
estimate the time in 6-31G(d).
\end{exercise}

\begin{exercise}[$\star\star$]\label{exo:b3:computational-chemistry:8}
% ledger: ie4:H2
The STO-3G orbital energy of \ce{H2} at \qty{1.4}{\bohr} is
\qty{-0.5782}{\hartree}. Give Koopmans' estimate of the ionisation energy in eV
and compare with \qty{15.43}{eV}. Name two reasons why they differ and their
directions.
\end{exercise}

\begin{exercise}[$\star\star$]\label{exo:b3:computational-chemistry:9}
In STO-3G, the energy of one hydrogen atom is \qty{-0.4666}{\hartree} and the RHF
energy of \ce{H2} at \qty{10}{\bohr} is \qty{-0.5960}{\hartree}. How far above two
atoms is the RHF energy at that distance, in eV? Explain with the ionic
terms.
\end{exercise}

\begin{exercise}[$\star\star\star$]\label{exo:b3:computational-chemistry:10}
For $\phi = c_1\chi_1 + c_2\chi_2$ with real functions, write $E(c_1,c_2)$ and
derive the two secular equations by setting $\partial E/\partial c_1 =
\partial E/\partial c_2 = 0$.
\end{exercise}

\begin{exercise}[$\star\star\star$]\label{exo:b3:computational-chemistry:11}
% ledger: aw:H, diat:H2.we
The RHF/STO-3G energy of \ce{H2} is $-1.116871$, $-1.117501$ and
$\qty{-1.116714}{\hartree}$ at $R = 1.30$, 1.35 and \qty{1.40}{\bohr}. Estimate
the \omterm{def:b3:quantum-model-systems:oscillator}{force constant} by a finite difference, convert it to \unit{N/m}
($\qty{1}{\hartree\per\bohr\squared} = \qty{1556.9}{N/m}$), and compute the
harmonic wavenumber. Compare with $\tilde\omega_e = \qty{4401}{cm^{-1}}$.
\end{exercise}

\begin{exercise}[$\star\star\star$]\label{exo:b3:computational-chemistry:12}
You must compute the energy difference between two conformers of a sugar,
known to differ by a few \unit{kJ/mol}, with an intramolecular hydrogen bond
in one of them. Choose and justify a method, a \omterm{def:b3:computational-chemistry:basis}{basis set} and the checks you
would run.
\end{exercise}

\section{Problem: HeH\texorpdfstring{$^+$}{+}, the First Molecule}

\begin{problem}[{Weekend problem --- a minimal-basis Hartree--Fock
calculation of the helium hydride ion: the basis, the core guess, the
converged orbital, and the binding of a proton to helium}]
\label{pb:b3:computational-chemistry:1}
% ledger: bse:He.sto3g.a1, bse:He.sto3g.a2, bse:He.sto3g.a3, bse:H.sto3g.a1, pa:He, ie:He, ie4:He2
\ce{HeH+} is computed at $R = \qty{1.4632}{\bohr}$ in the STO-3G basis (He
exponents 6.3624, 1.1589, 0.31365; H exponents 3.4253, 0.62391, 0.16886; the
same three coefficients). Function 1 is on He, function 2 on H. The program
gives (atomic units)
\[
\mathbf S = \begin{pmatrix}1 & 0.5368\\ 0.5368 & 1\end{pmatrix},\quad
\mathbf h = \begin{pmatrix}-2.5983 & -1.4318\\ -1.4318 & -1.7318\end{pmatrix},\quad
\mathbf F_{\mathrm{final}} = \begin{pmatrix}-1.5902 & -1.0610\\ -1.0610 & -0.8340
\end{pmatrix},
\]
the occupied orbital $\phi = 0.8766\chi_1 + 0.2025\chi_2$, the orbital energies
$-1.6328$ and \qty{-0.1725}{\hartree}, and the SCF energies $-2.7978$, $-2.8404$,
$-2.8418$, $-2.8418$. A helium atom in the same basis has
\qty{-2.8078}{\hartree}. $\qty{1}{\hartree} = \qty{27.211}{eV} =
\qty{2625.5}{kJ/mol}$.

\medskip\noindent\textbf{Part I --- The basis.}
\begin{enumerate}
  \item What does ``STO-3G'' mean?
  \item Why are the helium exponents larger than the hydrogen ones? Check
        that their ratio is $(1.69/1.24)^2$.
  \item How many basis functions, electrons and occupied orbitals are there?
  \item Compute the nuclear repulsion energy.
  \item Interpret $S_{12} = 0.5368$.
  \item Why is $h_{11}$ lower than $h_{22}$?
\end{enumerate}

\medskip\noindent\textbf{Part II --- The core guess.}
\begin{enumerate}[resume]
  \item Write the secular equation $\det(\mathbf h - \varepsilon\mathbf S) = 0$ as a
        quadratic in $\varepsilon$.
  \item Solve it.
  \item Why is this only a starting point?
  \item Give the formula of the Fock matrix in terms of $\mathbf h$, the
        density matrix and the two-electron integrals.
  \item Which electron interactions does $\mathbf F$ contain that $\mathbf h$ lacks?
  \item Explain why the procedure must be iterated.
\end{enumerate}

\medskip\noindent\textbf{Part III --- The converged orbital.}
\begin{enumerate}[resume]
  \item Check that $\phi$ is normalised.
  \item The populations are $P\mathbf S$ diagonal elements, $2c_1^2 + 2c_1c_2S$
        on He and $2c_2^2 + 2c_1c_2S$ on H. Compute them.
  \item Where does the positive charge sit? Is the molecule better described
        as \ce{He + H+} or \ce{He+ + H}?
  \item Use \omterm{thm:b3:computational-chemistry:koopmans}{Koopmans' theorem} to estimate the energy needed to remove an
        electron from \ce{HeH+}.
  \item Compute the electronic energy $E_{\mathrm{el}} = E - V_{\mathrm{nn}}$.
  \item How many cycles were needed to converge to \qty{e-4}{\hartree}? Why did
        the energy decrease at each cycle?
  \item With helium exponents scaled to $\zeta = 2.0925$ instead of 1.69, the
        same program gives \qty{-2.8607}{\hartree}. Which basis is better, and
        why can you say so?
  \item What does the virtual orbital energy, \qty{-0.1725}{\hartree}, represent?
\end{enumerate}

\medskip\noindent\textbf{Part IV --- The binding of a proton.}
\begin{enumerate}[resume]
  \item What is the energy of a bare proton? Of \ce{He + H+} far apart, in this
        basis?
  \item Compute the energy released by \ce{He + H+ -> HeH+} at this level, in
        \unit{kJ/mol}.
  \item The measured proton affinity of helium is \qty{177.8}{kJ/mol}.
        Comment.
  \item The exact energy of helium is minus the sum of its ionisation
        energies, 24.5874 and \qty{54.4178}{eV}. Compute it in \unit{\hartree} and
        the error of the STO-3G helium atom.
  \item Why can an error of this size still give useful molecular
        geometries?
  \item State the result: the RHF/STO-3G total energy of \ce{HeH+} at
        \qty{1.4632}{\bohr}.
\end{enumerate}
\end{problem}
